% \section{S-polynomial construction}

% To guarantee the applicability of the proposed algorithm, we must ensure that, given a number of variables \(n\), a degree \(d\), and a polynomial set of cardinality \(\binom{n+d-1}{d}\) whose leading monomials exhaust all monomials of total degree \(d\), a sufficient number of \(S\)-polynomials can be constructed (with respect to \Cref{condition-on-S-poly}). 
% In particular, we need to generate \(\binom{n+d-2}{d-1}\) linearly independent polynomials.

% %\paragraph{Dividing the monomials of total degree \(d\) w.r.t their higher degree component}
% \paragraph*{Partition of degree-\(d\) monomials by maximal exponent (highest degree component)}

% In accordance with \Cref{condition-on-S-poly}, \(S\)-polynomials are constructed by pairing polynomials whose leading monomials either have the same highest exponent or differ by one, for \(1 \le i \le d\).

% We therefore partition the set of monomials of total degree \(d\) according to the maximal exponent appearing in the monomial. 
% Let $\mathcal{S}_i :=\{\alpha\in\mathbb{N}^n : |\alpha|=d,\ \max(\alpha_i) _{1\le i\le n} = i \}$ denote the set of degree-\(d\) monomials whose maximal exponent is equal to \(i\).

% For \(\lceil d/2 \rceil \le i \le d\), the cardinality of \(\mathcal{S}_i\) is given by
% \begin{equation}
%     |\mathcal{S}_i| = \binom{n}{1}\binom{(n-1) + (d-i) -1}{(d-i)}.
% \end{equation} 
% % TODO: check for parity of d/2

% For \(1 \le i < \lceil d/2 \rceil\), the size of \(\mathcal{S}_i\) is obtained by imposing an upper bound \(i\) on the maximum exponent of each variable. 
% An inclusion-exclusion argument yields
% \begin{equation}
%     |\mathcal{S}_i| =
%     \sum_{j=0}^{\lfloor d/(i+1) \rfloor}
%     \binom{n}{j}
%     \binom{n + d - j(i+1) - 1}{d - j(i+1)}.
% \end{equation}

% By construction, the sets \(\{\mathcal{S}_i\} _{i=1}^{d}\) form a partition of the degree-\(d\) monomials, and therefore
% \[
%     \sum_{i=1}^{d} |\mathcal{S}_i| = \binom{n+d-1}{d}.
% \]

% Using this partition, \Cref{alg:pairs-construction} determines the number of \(S\)-polynomials that can be generated by pairing polynomials belonging either to the same set or to consecutive sets.

% % \begin{algorithm}[t!]
% %     \caption{$S$-polynomial Construction Check}\label{alg:pairs-construction}
% %     \begin{algorithmic}[1]
% %         \Statex{\hspace*{-\algorithmicindent} \textbf{Input:} $n \ge 1$: dimension, $1\le i\le d$: current degree}
% %         \Statex{\hspace*{-\algorithmicindent} \textbf{Output:} \texttt{True} if at least $\binom{n+i-2}{i-1}$ $S$-polynomial of the form~\Cref{eq:spoly} can be constructed, \texttt{False} otherwise}
% %         \Statex{}
% %         \Statex{\texttt{firstHalf = 0}}
% %         \Statex{\texttt{secondHalf = 0}}
% %         \Statex{\texttt{if d\% 2 == 0 then half = d//2+1 else half = float(d)/2}}
% %         \For{\texttt{k = d,\ldots, ceil (half)-1} by -1}
% %             \For{\texttt{j = 0,\ldots, ceil (half)-1 by -1}}
% %                 \texttt{total +=} $\binom{n}{j}\binom{n + (d-k) -1}{n-1}$
% %             \EndFor{}
% %         \EndFor{}
% %     \end{algorithmic}
% % \end{algorithm}



% \begin{algorithm}
% \caption{Gray-code composition generator: \texttt{gray\_comp}(n,d)}
% \label{alg:gray_comp}
% \begin{algorithmic}[1]
% \If{$P$ is \textbf{null}$ $}
%   \State $P\leftarrow []$
% \EndIf
% \If{$n = 1$}
%   \Return $P\mathbin{\Vert}[d]$
% \EndIf
% \For{$i=0$ to $d$}
%   \If{$i\bmod 2 = 0$}
%     \State \textit{/* even index: recurse and yield in natural order */}
%     \State yield all elements of \texttt{gray\_comp}($n-1,\;d-i,\;P\mathbin{\Vert}[i]$)
%   \Else
%     \State \textit{/* odd index: collect, reverse, then yield */}
%     \State $B \leftarrow$ list(\texttt{gray\_comp}($n-1,\;d-i,\;P\mathbin{\Vert}[i]$))
%     \State reverse($B$)
%     \State yield all elements of $B$
%   \EndIf
% \EndFor
% \end{algorithmic}
% \end{algorithm}

% \begin{algorithm}
% \caption{Generate admissible pairs: \texttt{gen\_admissible\_pair}(n,d)}
% \label{alg:gen_admissible_pair}
% \begin{algorithmic}[1]
% \State $pairs \leftarrow []$
% \State $to\_couple \leftarrow 0$
% \ForAll{$a$ in \texttt{gray\_comp}($n,d$)}
%   \If{$|pairs| = \binom{n+d-2}{d-1}$}
%     \State \textbf{break}
%   \EndIf
%   \If{$to\_couple \bmod 2 = 0$}
%     \State append $[a]$ to $pairs$
%   \Else
%     \State append $a$ to the last element of $pairs$
%   \EndIf
%   \State $to\_couple \leftarrow to\_couple + 1$
% \EndFor
% \State \Return $pairs$
% \end{algorithmic}
% \end{algorithm}



% % \begin{lemma}\label{lemma:number_spoly}
% % Given a degree \(d\), a number of variables \(n\), and a polynomial set of size \(\binom{n+d-1}{d}\) whose leading monomials exhaust all monomials of degree \(d\), \Cref{alg:pairs-construction} generates at least \(\binom{n+d-2}{d-1}\) \(S\)-polynomials if and only if \(d \le n-1\).
% % \end{lemma}



% % \begin{proof}
% %     To generate at least \(\binom{n+d-2}{d-1}\) $S$-polynomials we need at least \(2\binom{n+d-2}{d-1}\). 
% %     It is therefore necessary that
% %     \[
% %     \begin{split}
% %        \binom{n + d - 1}{d} &\ge 2\binom{n+d-2}{d-1} \implies \frac{(n+d-1)\cdots(n+1)(n)}{d!} \ge 2\frac{(n+d-2)\cdots(n)}{(d-1)!}\\
% %        \frac{(n+d-1)}{d} &\ge 2 \implies n+d-1 \ge 2d \implies  n-1 \ge d.
% %     \end{split}
% %     \]
% %     As a result, if \(d \le n-1\) we can generate enough $S$-polynomials. 
% % \end{proof}


% % The criterion described above can be implemented efficiently by organizing the polynomials in $\mathcal{D}_{i+1}$ according to the structure of their leading monomials.
% % Specifically, a polynomial $f \in \mathcal{D}_{i+1}$ with $\mathrm{LM}(f) = \mathbf{x}^\alpha$ is placed in the group indexed by $\max_j \alpha_j$.

% % During a $\deg_i$ iteration, we form pairs by selecting a polynomial from the group indexed by $i$ and pairing it with a polynomial from the group indexed by $i-1$. 
% % Once a polynomial has been used in a pair, it is marked as unavailable for subsequent pairings, i.e., it will not be considered for any other pair.  
% % This process is repeated until all polynomials in the group indexed by $i$ have been processed, or until no admissible pairs can be formed with the group indexed by $i-1$.  

% % When group $i$ runs out of polynomials, the iteration proceeds with the next pairs of groups, namely $i-1$ and $i-2$. 
% % At each step, the number of pairs that can be formed is given by the smaller cardinality of the two groups.  
% % More generally, if at some point the group $j-1$ contains more polynomials than $j-2$, and all polynomials in $j-2$ have already been paired, we then consider the next available lower groups, $j-3$ and $j-4$, if they exist.
% % If additional pairs are required to reach the desired number of $S$-polynomials, we also allow forming pairs within the same group using polynomials that have not yet been paired.  


% % \begin{proof}
% % To generate at least \(\binom{n+d-2}{d-1}\) reduced \(S\)-polynomials, at least \(2\binom{n+d-2}{d-1}\) degree-\(d\) polynomials are required. It is therefore necessary that
% % \[
% %     \binom{n+d-1}{d} \ge 2\binom{n+d-2}{d-1}.
% % \]
% % Since
% % \[
% %     \frac{\binom{n+d-1}{d}}{\binom{n+d-2}{d-1}} = \frac{n+d-1}{d},
% % \]
% % the inequality holds if and only if \(n+d-1 \ge 2d\), that is, \(d \le n-1\).
% % \end{proof}

% %%%%%%%%%%%%%%%%%%%%%%%%%%%%%%%%%%%
% %%%%%%%%%%%%%%%%%%%%%%%%%%%%%%%%%%%


% % \begin{figure}[b!]
% % \centering
% % \resizebox{0.55\linewidth}{!}{
% %   \begin{tikzpicture}[
  matparens/.style={
    inner sep=3pt,
    left delimiter={(},
    right delimiter={)}
  },
  block/.style={draw=gray, thick, inner sep=2pt, minimum height=12mm, minimum width=12mm, font=\large\bfseries},
  bullet/.style={draw=gray, fill=gray!30, minimum width=12mm, minimum height=12mm},
  arrow/.style={->, thick, line cap=round},
  node distance=2mm
]

% -------------------
% BLOCK 1) Samples
% -------------------
\node[matparens] (A1) {
    \begin{tikzpicture}
        \node[gray, matrix={1}{1}, fill=lipicsBulletGray] (R) {};
        \draw[gray, line width=0.70mm] (R.north west) -- (R.south east) -- cycle;
        \node[font=\footnotesize, fill=white, draw=gray, line width=0.4pt, inner sep=1pt] at (R.center) {$i+1$};

        \node[gray, matrix={1}{3}, fill=gray, right=1mm of R] (T) {};
        \node[black, align=center, text width=21mm, anchor=center] at (T.center) {degree $<i+1$};

        \node[draw, dotted, thick, rounded corners, inner sep=3pt, fit=(R)(T)] (boxd) {};

    \end{tikzpicture}
  };


% -------------------
% 2) DiagA
% -------------------
% \node[matparens, below=of A1] (A2) {
%     \begin{tikzpicture}
%         \node[lipicsGray, matrix={1}{1}, fill=white] (P) {};
%         \begin{scope}
%         \clip (P.north west) -- (P.south east) -- (P.north east) -- cycle;
%         \fill[lipicsBulletGray] (P.north west) rectangle (P.south east);
%         \end{scope}
%         \draw[lipicsGray, line width=0.70mm] (P.north west) -- (P.south east);
%         \draw[lipicsGray, line width=0.4pt] (P.north west) -- (P.north east) -- (P.south east) -- (P.south west) -- cycle;

%         \node[font=\large, fill=white, draw=lipicsGray, line width=0.4pt, inner sep=2pt, anchor=center] at (P.center) {\large$i$};

%         \node[lipicsBulletGray, matrix={1}{3}, fill=lipicsLightGray, right=2mm of P] (L) {};
%         \node[black, align=center, text width=16mm, anchor=center] at (L.center) {degree $<i$};
%     \end{tikzpicture}
% };

% % Arrow A1 -> A2
% \draw[arrow] (A1.south) -- (A2.north) node[midway, right=1mm, font=\scriptsize, align=center] {\texttt{DiagA}};

% -------------------
% BLOCK 2) DiagA + DiagB
% -------------------
\node[matparens, below=of A1] (B2) {
    \begin{tikzpicture}
        \node[gray, matrix={1}{1}, fill=white] (P) {};
        \draw[gray, line width=0.70mm] (P.north west) -- (P.south east) -- cycle;
        \node[font=\footnotesize, fill=white, draw=gray, line width=0.4pt, inner sep=1pt, anchor=center] at (P.center) {$i+1$};

        \node[gray, matrix={1}{3}, fill=gray, right=1mm of P] (L) {};
        \node[black, align=center, text width=21mm, anchor=center] at (L.center) {degree $<i+1$};

        \node[draw, dotted, thick, rounded corners, inner sep=3pt, fit=(P)(L)] (boxd) {};
    \end{tikzpicture}
};

% Arrow A2 -> B2
\draw[arrow]
  (A1.south) -- (B2.north)
  node[midway, right=1mm, font=\scriptsize, align=center] {\texttt{DiagA + DiagB}};


% -------------------
% BLOCK 3) Spoly generation
% -------------------
\node[matparens, below=of B2] (A1b) {
\begin{tikzpicture}
    \node[gray, matrix={0.7}{1}, fill=gray] (K) {};
    \node[font=\footnotesize, fill=white, draw=lipicsGray, line width=0.4pt, inner sep=1pt, anchor=center] at (K.center) {$i+1$};
    
    \node[gray, matrix={0.7}{0.7}, fill=gray, right=1mm of K] (R) {};
    \draw[gray, line width=0.70mm] (R.north west) -- (R.south east) -- cycle;
    \node[font=\footnotesize, fill=white, draw=lipicsGray, line width=0.4pt, inner sep=1pt, anchor=center] at (R.center) {$i$};


    \node[gray, matrix={0.7}{2.2}, fill=gray, right=1mm of R] (T) {};
    \node[black, align=center, text width=16mm, anchor=center, font=\footnotesize] at (T.center) {degree $<i$};

    \node[draw, dotted, thick, rounded corners, inner sep=3pt, fit=(K)(R)(T)] (boxd) {};
\end{tikzpicture}
};

% Arrow B2 -> A1b
\draw[arrow] (B2.south) -- (A1b.north) node[midway, right=1mm, font=\scriptsize, align=center] {\texttt{Spoly} (gen.)};

% -------------------
% BLOCK 4) Spoly reduction
% -------------------
\node[matparens, below=of A1b] (B3) {
\begin{tikzpicture}
    %\node[draw=lipicsGray, dashed, line width=0.6pt, fill=none, matrix={0.7}{1}] (K) {};
    \node[matrix={0.7}{1}, draw=none, fill=none] (K) {};

    \node[gray, matrix={0.7}{0.7}, fill=gray, right=1mm of K] (R) {};
    \draw[gray, line width=0.70mm] (R.north west) -- (R.south east) -- cycle;
    \node[font=\footnotesize, fill=white, draw=gray, line width=0.4pt, inner sep=1pt, anchor=center] at (R.center) {$i$};

    \node[gray, matrix={0.7}{2.2}, fill=gray, right=1mm of R] (T) {};
    \node[black, align=center, text width=16mm, anchor=center, font=\footnotesize] at (T.center) {degree $<i$};

    \node[draw, dotted, thick, rounded corners, inner sep=3pt, fit=(R)(T)] (boxd) {};
\end{tikzpicture}
};

% Arrow B2 -> A1b
\draw[arrow] (A1b.south) -- (B3.north) node[midway, right=1mm, font=\scriptsize, align=center] {\texttt{Spoly} (red.)};

% -------------------
% BLOCK 5) Diagonalization
% -------------------
\node[matparens, below=of B3] (AA1) {
    \begin{tikzpicture}
        % \node[draw=lipicsGray, dashed, line width=0.6pt, fill=none, matrix={0.7}{1}] (K) {};
        \node[matrix={0.7}{1}, draw=none, fill=none] (K) {};

        \node[gray, matrix={0.7}{0.7}, fill=white, right=1mm of K] (P) {};
        \draw[gray, line width=0.70mm] (P.north west) -- (P.south east) -- cycle;
        \node[font=\footnotesize, fill=white, draw=gray, line width=0.4pt, inner sep=1pt, anchor=center] at (P.center) {$i$};

        \node[gray, matrix={0.7}{2.2}, fill=gray, right=1mm of P] (L) {};
        \node[font=\footnotesize, black, align=center, text width=16mm, anchor=center] at (L.center) {degree $<i$};

        \node[draw, dotted, thick, rounded corners, inner sep=3pt, fit=(P)(L)] (boxd) {};
    \end{tikzpicture}
};

% Arrow A2 -> B2
\draw[arrow]
  (B3.south) -- (AA1.north)
  node[midway, right=1mm, font=\scriptsize, align=center] {\texttt{DiagA + DiagB}};


% Arrow B3 -> A1
% \draw[arrow, dashed]
%   ($(AA1.east)+(3mm,0)$)
%   -- ($(AA1.east)+(7mm,0)$)
%   -- ($(B2.east)+(7mm,0)$)
%   -- ($(B2.east)+(3mm,0)$);

\end{tikzpicture}
% % }
% % \caption{Not decided.}\label{fig:oooop}
% % \end{figure}




